\documentclass[11pt]{article}
\usepackage{amsmath,amssymb,amsthm,geometry,hyperref}
\geometry{margin=1in}
\newtheorem{lemma}{Lemma}
\newtheorem{theorem}{Conditional theorem}
\newcommand{\R}{\mathbb R}
\newcommand{\norm}[1]{\left\|#1\right\|}
\begin{document}
\title{A finite-accuracy kernel criterion for positive-energy nontriviality}
\date{16 October 2026}\maketitle
\noindent Status: the free infinite-lattice calculation and elementary
comparison lemmas below are derived here. Their application to the deposited
Yang--Mills witness requires the specifically stated fluctuation-curvature
estimate and the interacting witness itself. Neither is certified here.

\section{A free lattice calculation independent of the missing companion note}

Use unit lattice coordinates initially. Put
$s_\mu(q)=2\sin(q_\mu/2)$ and $\lambda(q)=\sum_\mu s_\mu(q)^2$.
In midpoint conventions for links and plaquettes the Fourier symbol of the
covariance of one colour of the exact two-form Gaussian is
\[
m_{\mu\nu,\rho\sigma}(q)=
\frac{\delta_{\nu\sigma}s_\mu s_\rho-\delta_{\nu\rho}s_\mu s_\sigma
-\delta_{\mu\sigma}s_\nu s_\rho+\delta_{\mu\rho}s_\nu s_\sigma}
{\lambda(q)}.
\]
The indices satisfy $\mu<\nu,\rho<\sigma$. Equivalently this is the
projection $s\wedge\iota_s/|s|^2$ on exact two-forms.
The symbol is interpreted on staggered lattice cosets; multiplying by
the fixed midpoint phase gives an ordinary periodic lattice multiplier.
Let $m_0(q)$ be the same homogeneous expression with $s_\mu$ replaced by
$q_\mu$.

\begin{lemma}[Off-diagonal free covariance]
Let $D_{PQ}(r)$ be the continuum Maxwell curvature covariance with
potential covariance $\delta_{\mu\nu}/(4\pi^2|r|^2)$.
For bare spacing $a$, and distinct plaquette midpoints with
$|r|=|x_p-x_q|\geq C_0a$,
\[
\left|P_{{\rm ex},a}(p,q)-a^4D_{\pi(p)\pi(q)}(r)\right|
\leq C a^6|r|^{-6}.
\tag{1}
\]
Here $C,C_0$ are finite numerical constants for this lattice discretization.
\end{lemma}
\begin{proof}
In a small ball about zero, Taylor expansion of sine and differentiation
of the quotient give, for every multi-index $\alpha$,
\[
|\partial^\alpha(m-m_0)(q)|\leq C_\alpha|q|^{2-|\alpha|}.
\]
Choose a smooth cutoff $\chi$ supported in that ball and equal to one
near zero. A dyadic partition into $|q|\asymp2^{-j}$ writes
$\chi(m-m_0)=\sum_{j\geq0}r_j$. Integrating by parts $M>6$ times on
each compact annulus, or using its rescaled smooth Fourier transform,
gives
\[
|\widehat r_j(z)|\leq C_M2^{-6j}
(1+2^{-j}|z|)^{-M}.
\]
Summing above and below $2^{-j}|z|\asymp1$ gives
$|\sum_j\widehat r_j(z)|\leq C|z|^{-6}$ for $|z|\geq1$.
For precision, if $n$ is the integer basepoint difference and
$\delta_P,\delta_Q$ are the two orientation midpoint offsets, use the
same factor $e^{iq\cdot(\delta_P-\delta_Q)}$ on both $m$ and $m_0$.
Their difference remains $O(|q|^2)$ with the stated derivative bounds,
and its Fourier coefficient at $n$ compares the two kernels at
$z=n+\delta_P-\delta_Q=r/a$. Applying the phase to only one symbol
would introduce a spurious first-order remainder.

The lattice multiplier away from zero is smooth and periodic in the
basepoint convention, so its Fourier coefficients decrease faster than
any prescribed inverse power. The continuum distribution with
multiplier $m_0$ is $D$ away from zero. Removing its small-frequency
cutoff produces a term decreasing faster than any inverse power at
large nonzero $z$: all sufficiently high derivatives of
$(1-\chi)m_0$ are integrable, so repeated integration by parts in the
tempered Fourier transform applies there. Contact terms are irrelevant
at these separated midpoints.
Consequently the unit-lattice error is $O(|z|^{-6})$.
Set $z=r/a$ and use homogeneity $D(r/a)=a^4D(r)$.
\end{proof}

This proof establishes the infinite-lattice estimate. It does not silently
assert a quantitative torus comparison. The finite-torus free covariance
converges to this one at fixed spacing and fixed plaquettes by its Fourier
Riemann sum, since the multiplier is bounded and discontinuous only at
one momentum. For a block-conditioned covariance, existence and identity
of the corresponding volume limit must also be justified.

\section{The exact remaining free-background input}

Suppose the background covariance is
$C_{{\rm bg},a,h}=P_{{\rm ex},a}-C_{{\rm fl},a,h}$.
Here $h$ is the physical witness-cell scale, not the bare spacing.
The sufficient input is a \emph{curvature} estimate
\[
|C_{{\rm fl},a,h}(p,q)|
\leq A a^4h^{-4}\exp(-c|x_p-x_q|/h),
\qquad |x_p-x_q|\geq h,
\tag{2}
\]
with $A,c>0$ independent of $a,h$ and full volume.
A bound on an undifferentiated potential propagator is insufficient
without the correctly normalized discrete curl estimates.

One possible way of obtaining (2) is a scale decomposition with
curvature terms bounded by
$A_0a^4h_j^{-4}e^{-c_0r/h_j}$, $h_j=hL^{-j}$.
Indeed for $r/h\geq1$,
\[
\sum_{j\geq0}L^{4j}e^{-c_0(r/h)L^j}
\leq C_{L,c_0}e^{-c_0r/(2h)}.
\]
Factor out $e^{-c_0r/(2h)}$ and dominate the remaining series by
$\sum_jL^{4j}e^{-c_0L^j/2}<\infty$.
This summation is elementary; the scale decomposition and its
curl estimates are the model/operator identification obligation.

\section{Only fixed relative accuracy is necessary}

Let $f_d$ be the bump used in Notebook 01, centered at $3de_0/4$
and of radius $d/4$. Denote its reflected near pair by $-$ and its
reflected pair with time translation $4d$ by $+$.
Define
\[
K_\pm=\frac6{\pi^4}
\iint f_\pm(x)g_\pm(y)|x-y|^{-8}\,dx\,dy.
\]
They are positive and $K_+/K_-\leq(2/5)^8$.
Both are scale independent under sup-normalized dilation.
Use the same normalization for the discrete background kernel:
\[
 K_{{\rm bg},\pm}=\sum_{p,q}
 f_\pm(x_p)g_\pm(x_q)
 |C_{{\rm bg},a,h}(p,q)|^2,
\]
where each positively oriented plaquette occurs once. The covariance is
for one colour. A common energy-density or colour normalization may
multiply both this expression and $K_\pm$ and cancels from every ratio.
The factors $a^4$ in each covariance provide the $a^8$ Riemann-sum
weight; no additional lattice-volume factors belong in this definition.

\begin{lemma}[Finite-accuracy criterion]
If the corresponding background quadratic kernels satisfy
\[
(1-\epsilon)K_\pm\leq K_{{\rm bg},\pm}
\leq(1+\epsilon)K_\pm,\qquad 0\leq\epsilon<1,
\tag{3}
\]
and the interacting witness has a common positive factor $b g_d^4$
and prefactor ratio at most twelve for these two pairs, then
\[
\frac{C_d(4d)}{C_d(0)}
\leq12\frac{1+\epsilon}{1-\epsilon}(2/5)^8.
\tag{4}
\]
In particular $\epsilon=1/3$ gives $0.01572864<0.02$.
\end{lemma}
\begin{proof}
Use the upper witness on the far pair, the lower witness on the near
pair, and cancel the common running coupling. Substitute (3).
\end{proof}

Thus a one-third relative kernel comparison suffices. The whole
asymptotic expansion, its exact error coefficient, and convergence
for arbitrary shapes are unnecessary for this particular witness.

\begin{theorem}[A sufficient two-scale condition]
Assume (1), (2), and the stated infinite-volume background covariance.
For the two fixed shapes there is $C_*>0$, depending on their shapes
and $A,c$ but not on $a,h,d$, such that (3) holds with
\[
\epsilon\leq C_*\left\{\frac ad+
\left(\frac dh\right)^4e^{-c\,d/h}\right\}
\tag{5}
\]
whenever $a/d$ is sufficiently small and $d/h\geq1$, with a harmless
fixed reduction in $c$ if needed.
Consequently finite thresholds $a/d\leq\eta_*$ and $d/h\geq R_*$
ensure $\epsilon\leq1/3$.
\end{theorem}
\begin{proof}
All distances within either pair are in $[d,6d]$.
Equation (1) gives entrywise relative error $O((a/d)^2)$ against
$a^4d^{-4}$, and (2) gives
$O((d/h)^4e^{-c\,d/h})$ against the same scale.
Group the six orientations at each lattice basepoint into a $6$ by $6$
matrix. Moving their midpoints to a common basepoint changes the
continuum matrix by $O(a/d)$ relative to $d^{-4}$, because
$|\nabla D(r)|\leq C|r|^{-5}$.
For such matrices,
\[
|\norm{D+E}_{\rm F}^2-\norm D_{\rm F}^2|
\leq2\norm D_{\rm F}\norm E_{\rm F}+\norm E_{\rm F}^2,
\qquad
\norm D(r)_{\rm F}^2=\frac6{\pi^4|r|^8}.
\]
The weights are bounded, compactly supported fixed Lipschitz shapes
after scaling by $d$. Their midpoint shifts and Riemann sums cost
$O(a/d)$ in the rescaled integral. The positive continuum denominators
$K_\pm$ are fixed nonzero constants for these shapes.
Dividing the absolute error by them yields (5), absorbing quadratic
errors for sufficiently small parameters into its constant.
The exponential dominates the fourth power as $d/h\to\infty$.
\end{proof}

The deposited witness chooses $h=L^{-s(d)}$ with
$d/h$ growing polylogarithmically as $d\downarrow0$.
This supplies the requisite scale separation \emph{if} (2) is verified
for that exact block map. First take the cutoff fine at fixed $d$,
then choose $d$ small; no claim of arbitrary joint-limit uniformity is
needed for the existence of one such probe.

\section{Spectral consequence and boundary of the result}

If a limiting Hilbert space with $H\geq0$ has been reconstructed,
(4) with $q=0.02$ gives
\[
\norm{\mathbf1_{(1/(4d),\infty)}(H)v_d}^2
\geq(1-0.02e)\norm{v_d}^2>0.
\]
This follows by integrating $e^{-4dE}\geq e^{-1}$ on
$E\leq1/(4d)$. The vector $(1-e^{-4dH})v_d$ is nonzero and
orthogonal to $\ker H$.

Remaining obligations are exact block-map/covariance identification,
the uniform curvature estimate (2), the interacting lower/upper
witness, and a valid reconstruction. The calculation does not prove
those inputs, uniqueness, a mass gap, or Yang--Mills identification.
\end{document}
